\begin{frame}{역사: ``신뢰 영역''은 50년 된 아이디어}
  \small
  \begin{tabular}{@{}ll@{}}
    \toprule
    시기 & 사건 \\
    \midrule
    1960\,$\sim$\,70년대 & \kb{미분없는 최적화}의 고전 기법:
      ``목표함수의 2차 근사가 잘 맞는 작은 영역 안에서만
      근사함수를 최적화한다'' \\
    1971 & Maritz — 이 분리를 \textbf{확률적 시뮬레이션} 문맥에서
      처리 \\
    1981 & Moore — 미분없는 최적화에서 정식화 \\
    2015 & \kb{TRPO} — ``KL 발산이 $\delta$ 안이면 policy
      gradient의 2차 근사가 안전하다''를 증명, 2차 계획으로
      매 스텝 해석 (보장은 아름답고, 구현은 무거움) \\
    2017 & \kb{PPO}(Schulman 외) — ``보장을 1차 기법(클리핑)으로
      근사해도 실전에서 거의 같은 안정성''이라는 관찰 위에 \\
    \bottomrule
  \end{tabular}
  \vskip0.8em
  \begin{block}{핵심 교훈}
    ``TRPO는 이론적으로 정답, PPO는 실용적으로 정답''이 아니라 —
    \kb{같은 신뢰 영역 원리를, 풀이를 쉬운 쪽으로 바꿨을 뿐}.
    근사된 목적함수를 멀리서 믿지 말고, 가까운 영역에서만
    믿는다는 태도 자체가 통계·최적화에서 오래된 지혜.
  \end{block}
\end{frame}
